\section{Methodology}
\label{sec:methodology}

This study adopts a \gls{slr} methodology following the \gls{prisma} 2020 guidelines to provide a transparent, systematic, and reproducible review of the existing literature on \gls{dectnr}. Unlike conventional narrative surveys, the \gls{prisma} framework employs a predefined review protocol that minimizes selection bias through a structured process comprising literature identification, screening, eligibility assessment, quality assessment, data extraction, and qualitative synthesis.

The review protocol was established prior to the literature search to ensure methodological consistency and reproducibility throughout the review process. As illustrated in Figure~\ref{fig:prisma}, the review consisted of nine sequential stages:
\begin{figure*}[!t]
\centering
\resizebox{\textwidth}{!}{%
\begin{tikzpicture}[
    font=\sffamily,
    phase/.style={
        rounded corners=4pt,
        draw=#1,
        very thick,
        fill=#1!4,
        minimum width=6.7cm,
        minimum height=17.8cm
    },
    header/.style={
        rounded corners=4pt,
        draw=#1,
        fill=#1,
        text=white,
        font=\bfseries\Large,
        minimum width=6.6cm,
        minimum height=1.55cm,
        align=center
    },
    purpose/.style={
        font=\bfseries\itshape,
        text width=5.55cm,
        align=center,
        minimum height=1.5cm
    },
    step/.style={
        rounded corners=3pt,
        draw=#1,
        fill=#1!7,
        line width=0.8pt,
        text width=5.45cm,
        minimum height=1.9cm,
        inner sep=5pt,
        align=left
    },
    output/.style={
        rounded corners=3pt,
        draw=#1,
        fill=#1!5,
        text width=5.45cm,
        minimum height=3.0cm,
        inner sep=5pt,
        align=left
    },
    flow/.style={
        -{Latex[length=2.4mm]},
        line width=0.9pt
    }
]

\definecolor{identblue}{RGB}{39,84,138}
\definecolor{screengreen}{RGB}{52,112,64}
\definecolor{eliggold}{RGB}{196,143,0}
\definecolor{synthpurple}{RGB}{101,60,140}

% ============================================================
% PHASE PANELS
% ============================================================

\node[phase=identblue] (p1) at (0,0) {};
\node[phase=screengreen, right=0.55cm of p1] (p2) {};
\node[phase=eliggold, right=0.55cm of p2] (p3) {};
\node[phase=synthpurple, right=0.55cm of p3] (p4) {};

% ============================================================
% HEADERS
% ============================================================

\node[header=identblue, anchor=north]
    (h1) at ([yshift=-0.15cm]p1.north)
    {1. IDENTIFICATION};

\node[header=screengreen, anchor=north]
    (h2) at ([yshift=-0.15cm]p2.north)
    {2. SCREENING};

\node[header=eliggold, anchor=north]
    (h3) at ([yshift=-0.15cm]p3.north)
    {3. ELIGIBILITY};

\node[header=synthpurple, anchor=north]
    (h4) at ([yshift=-0.15cm]p4.north)
    {4. INCLUDED STUDIES\\AND SYNTHESIS};

% ============================================================
% PURPOSE STATEMENTS
% ============================================================

\node[purpose, below=0.30cm of h1]
    (pu1)
    {Finding and collecting\\potentially relevant studies};

\node[purpose, below=0.30cm of h2]
    (pu2)
    {Filtering studies based on\\titles and abstracts};

\node[purpose, below=0.30cm of h3]
    (pu3)
    {Assessing full-text studies\\for eligibility and quality};

\node[purpose, below=0.30cm of h4]
    (pu4)
    {Extracting data and\\synthesizing evidence};

% ============================================================
% IDENTIFICATION
% ============================================================

\node[step=identblue, below=0.35cm of pu1] (id1) {
\textbf{1.1 Formulate research questions}\\
Define the scope, objectives, and research questions.
};

\node[step=identblue, below=0.22cm of id1] (id2) {
\textbf{1.2 Select databases}\\
IEEE Xplore, Scopus, ACM Digital Library, SpringerLink,
ScienceDirect, and Google Scholar.
};

\node[step=identblue, below=0.22cm of id2] (id3) {
\textbf{1.3 Develop search strategy}\\
Define keywords, synonyms, Boolean operators, and search strings.
};

\node[step=identblue, below=0.22cm of id3] (id4) {
\textbf{1.4 Search literature}\\
Run searches in the selected databases and export all records.
};

\node[step=identblue, below=0.22cm of id4] (id5) {
\textbf{1.5 Remove duplicates}\\
Merge all records and remove duplicate publications.
};

\node[output=identblue, anchor=south]
    (out1) at ([yshift=0.30cm]p1.south) {
\textbf{Key outputs}\\
$\bullet$ Search strategy\\
$\bullet$ Records retrieved $(n=\ldots)$\\
$\bullet$ Records after deduplication $(n=\ldots)$
};

\draw[flow] (id1) -- (id2);
\draw[flow] (id2) -- (id3);
\draw[flow] (id3) -- (id4);
\draw[flow] (id4) -- (id5);
\draw[flow] (id5) -- (out1);

% ============================================================
% SCREENING
% ============================================================

\node[step=screengreen, below=0.50cm of pu2] (sc1) {
\textbf{2.1 Title screening}\\
Exclude records that are clearly outside the review scope.
};

\node[step=screengreen, below=0.35cm of sc1] (sc2) {
\textbf{2.2 Abstract screening}\\
Apply inclusion and exclusion criteria to the abstracts.
};

\node[step=screengreen, below=0.35cm of sc2] (sc3) {
\textbf{2.3 Record exclusion reasons}\\
Document the principal reasons for excluding records.
};

\node[output=screengreen, anchor=south]
    (out2) at ([yshift=0.30cm]p2.south) {
\textbf{Key outputs}\\
$\bullet$ Records screened $(n=\ldots)$\\
$\bullet$ Records excluded $(n=\ldots)$\\
$\bullet$ Reports sought $(n=\ldots)$
};

\draw[flow] (sc1) -- (sc2);
\draw[flow] (sc2) -- (sc3);
\draw[flow] (sc3) -- (out2);

% ============================================================
% ELIGIBILITY
% ============================================================

\node[step=eliggold, below=0.50cm of pu3] (el1) {
\textbf{3.1 Retrieve full texts}\\
Obtain the complete publications for potentially eligible studies.
};

\node[step=eliggold, below=0.28cm of el1] (el2) {
\textbf{3.2 Assess full-text eligibility}\\
Apply the predefined inclusion and exclusion criteria.
};

\node[step=eliggold, below=0.28cm of el2] (el3) {
\textbf{3.3 Quality assessment}\\
Evaluate methodological quality using the quality checklist.
};

\node[step=eliggold, below=0.28cm of el3] (el4) {
\textbf{3.4 Record exclusion reasons}\\
Document reasons for exclusion after full-text assessment.
};

\node[output=eliggold, anchor=south]
    (out3) at ([yshift=0.30cm]p3.south) {
\textbf{Key outputs}\\
$\bullet$ Full texts assessed $(n=\ldots)$\\
$\bullet$ Full texts excluded $(n=\ldots)$\\
$\bullet$ Eligible studies $(n=\ldots)$
};

\draw[flow] (el1) -- (el2);
\draw[flow] (el2) -- (el3);
\draw[flow] (el3) -- (el4);
\draw[flow] (el4) -- (out3);

% ============================================================
% INCLUDED STUDIES AND SYNTHESIS
% ============================================================

\node[step=synthpurple, below=0.65cm of pu4] (in1) {
\textbf{4.1 Data extraction}\\
Extract predefined variables using a structured data-extraction form.
};

\node[step=synthpurple, below=0.45cm of in1] (in2) {
\textbf{4.2 Data synthesis}\\
Perform descriptive, thematic, quantitative, and comparative analysis.
};

\node[step=synthpurple, below=0.45cm of in2] (in3) {
\textbf{4.3 Report results}\\
Present findings, taxonomy, research gaps, and future directions.
};

\node[output=synthpurple, anchor=south]
    (out4) at ([yshift=0.30cm]p4.south) {
\textbf{Key outputs}\\
$\bullet$ Included studies $(n=59)$\\
$\bullet$ Extracted dataset\\
$\bullet$ Taxonomy and synthesis\\
$\bullet$ Research roadmap
};

\draw[flow] (in1) -- (in2);
\draw[flow] (in2) -- (in3);
\draw[flow] (in3) -- (out4);

% ============================================================
% CONNECTIONS BETWEEN PHASES
% ============================================================

\draw[flow]
    ([yshift=1.0cm]p1.east)
    -- ([yshift=1.0cm]p2.west);

\draw[flow]
    ([yshift=1.0cm]p2.east)
    -- ([yshift=1.0cm]p3.west);

\draw[flow]
    ([yshift=1.0cm]p3.east)
    -- ([yshift=1.0cm]p4.west);

\end{tikzpicture}%
}
\caption{PRISMA-oriented systematic review workflow used in this study.
The first three phases describe study identification, screening, and
eligibility assessment, while the final phase covers extraction,
synthesis, and reporting of the included evidence.}
\label{fig:review_workflow}
\end{figure*}
\begin{enumerate}
    \item Formulation of the \glspl{rq};
    \item Selection of digital literature databases;
    \item Development of the search strategy;
    \item Systematic literature retrieval;
    \item Duplicate detection and removal;
    \item Study selection using predefined inclusion and exclusion criteria;
    \item \gls{qa} of the selected studies;
    \item Structured data extraction;
    \item Qualitative and comparative data synthesis.
\end{enumerate}

%%%%%%%%%%%%%%%%%%%%%%%%%%%%%%%%%%%%%%%%%%%%%%%%%%%%%%%%%%%%%%%%%%%%%%%%%%%%%%

\subsection{Review Protocol}

A review protocol was defined before commencing the literature search to ensure the transparency, consistency, and reproducibility of the entire review process. The protocol specifies the research objectives, search strategy, study selection criteria, quality assessment procedure, data extraction framework, and synthesis methodology adopted in this study. Although the protocol was not formally registered in an external repository, all methodological decisions were established before the study selection process began.

Table~\ref{tab:protocol} summarizes the review protocol adopted in this work.

\begin{table}[!t]
\caption{Summary of the Review Protocol}
\label{tab:protocol}
\centering
\begin{tabular}{p{3.1cm}p{4.4cm}}
\toprule
\textbf{Item} & \textbf{Description} \\
\midrule
Review Method & \gls{prisma} 2020 \gls{slr} \\
Research Domain & \gls{dectnr} \\
Publication Period & January 2020 -- July 2026 \\
Language & English \\
Primary Databases & IEEE Xplore, Scopus, ACM Digital Library, SpringerLink, ScienceDirect \\
Complementary Search & Google Scholar (Backward and Forward Snowballing) \\
Search Fields & Title, Abstract, and Author Keywords \\
Duplicate Management & Zotero Reference Manager \\
Study Selection & \gls{prisma} Screening Process \\
Quality Assessment & Five quality criteria (0--10) \\
Data Extraction & Structured extraction template \\
Data Analysis & Descriptive, comparative, and thematic synthesis \\
\bottomrule
\end{tabular}
\end{table}

%%%%%%%%%%%%%%%%%%%%%%%%%%%%%%%%%%%%%%%%%%%%%%%%%%%%%%%%%%%%%%%%%%%%%%%%%%%%%%

\subsection{Research Questions}

The primary objective of this \gls{slr} is to provide a comprehensive overview of the current state of \gls{dectnr} research, identify emerging research trends, evaluate existing implementation and validation approaches, and highlight future research opportunities. To achieve these objectives, the review addresses the following \glspl{rq}.

\begin{itemize}

\item[\textbf{RQ1}] How has research on \gls{dectnr} evolved in terms of publication trends, publication types, publishers, and application domains?

\item[\textbf{RQ2}] Which technical aspects of \gls{dectnr}, including the \gls{phy}, \gls{mac}, routing, resource allocation, positioning, security, and industrial applications, have been investigated?

\item[\textbf{RQ3}] Which evaluation methodologies, implementation platforms, software tools, and hardware platforms have been employed to validate \gls{dectnr} systems?

\item[\textbf{RQ4}] Which \glspl{kpi} have been used to evaluate \gls{dectnr}, and what performance characteristics have been reported?

\item[\textbf{RQ5}] What are the principal scientific contributions and technical limitations reported in the existing literature?

\item[\textbf{RQ6}] Which research gaps remain insufficiently addressed in the current body of \gls{dectnr} literature?

\item[\textbf{RQ7}] What future research directions are most promising for advancing \gls{dectnr} towards large-scale \gls{iiot} deployments?

\end{itemize}
\subsection{Eligibility Criteria}

The inclusion and exclusion criteria were defined before the screening
process. A study was included if it satisfied all of the following criteria:

\begin{itemize}
    \item it primarily investigated \gls{dectnr} or DECT NR+;
    \item it presented a technical, experimental, analytical, architectural,
    or application-oriented contribution;
    \item it was written in English;
    \item the full text was available;
    \item it was published between January 2020 and July 2026;
    \item it provided sufficient methodological or technical detail to support
    data extraction.
\end{itemize}

A study was excluded if it:

\begin{itemize}
    \item mentioned \gls{dectnr} only briefly or as background;
    \item primarily addressed another wireless technology;
    \item was a duplicate record;
    \item lacked sufficient technical content;
    \item was not available in English;
    \item was an editorial, presentation slide set, poster, or non-technical
    document.
\end{itemize}

Studies that were relevant to the broader context of industrial wireless
communications but did not primarily focus on \gls{dectnr} were classified
as contextual studies and were excluded from the quantitative synthesis.
\subsection{Quality Assessment}

Each included study was assessed using five quality criteria:

\begin{enumerate}
    \item clarity of the research objective;
    \item appropriateness of the methodology;
    \item strength of the validation;
    \item clarity of the discussion and limitations;
    \item reproducibility of the study.
\end{enumerate}

Each criterion was scored on a three-point scale:

\begin{itemize}
    \item 0: criterion not addressed;
    \item 1: criterion partially addressed;
    \item 2: criterion fully addressed.
\end{itemize}

The total quality score therefore ranged from 0 to 10. Studies scoring
8--10 were classified as high quality, studies scoring 5--7 as moderate
quality, and studies scoring below 5 as low quality.
\subsection{Data Extraction}

A structured data extraction form was used to collect information from each
included study. The extracted fields comprised:

\begin{itemize}
    \item bibliographic information;
    \item publication year and publication type;
    \item primary and secondary research contributions;
    \item technical topics;
    \item evaluation methodology;
    \item hardware and software platforms;
    \item network topology and deployment environment;
    \item frequency band and bandwidth;
    \item evaluated \glspl{kpi};
    \item application domain;
    \item key findings;
    \item scientific contribution;
    \item limitations;
    \item identified research gaps;
    \item proposed future work.
\end{itemize}

The extracted information was stored in a structured spreadsheet database.
Binary one-hot coding was used for multi-label attributes such as technical
topics, evaluation methods, application domains, and performance metrics.
\subsection{Study Classification and Coding}

Each study was assigned exactly one primary contribution category, while
multiple secondary technical topics were allowed. This distinction prevented
double counting in the analysis of primary research contributions while
preserving the multidisciplinary nature of the literature.

Multi-label technical topics included the \gls{phy}, \gls{mac},
routing and mesh networking, scheduling and resource allocation,
positioning, security, energy efficiency, coexistence, artificial
intelligence, standardization, industrial applications, \gls{urllc},
and \gls{mmtc}.

Evaluation methodologies were coded independently as simulation,
analytical or mathematical modeling, experimental or testbed validation,
and survey or overview. A paper could therefore belong to more than one
methodological category.
\subsection{Reliability and Consistency Checks}

To improve coding reliability, the extracted data were reviewed in multiple
stages. Evaluation methods, technical topics, performance metrics, hardware
platforms, and application domains were cross-checked across the master
database and the corresponding binary classification matrices.

Where inconsistencies were identified, the full paper, abstract, methodology,
results, and conclusion sections were re-examined. The final database was
then frozen and used as the single source for all quantitative analyses,
tables, figures, and percentages reported in this review.
\subsection{Data Synthesis and Analysis}

The included studies were analyzed using descriptive, comparative, and
thematic synthesis.

Descriptive analysis was used to examine publication trends, publication
types, research categories, evaluation methods, hardware platforms,
application domains, and performance metrics.

Comparative analysis was used to examine differences between simulation-based,
analytical, and experimental studies.

Thematic synthesis was used to group the reported scientific contributions,
limitations, research gaps, and future research directions into recurring
themes.

Because technical topics, evaluation methods, and performance metrics were
coded as multi-label attributes, the corresponding percentages may sum to
more than 100\%.
\subsection{Threats to Validity}

Several threats to validity should be acknowledged. First, relevant studies
may have been missed because of differences in database indexing or search
terminology. This risk was reduced by searching multiple databases and
performing backward and forward snowballing.

Second, the classification of research topics and contributions involves
researcher judgment. To reduce subjectivity, a predefined coding framework
was applied consistently, and ambiguous cases were re-examined using the
full text.

Third, the rapidly evolving nature of \gls{dectnr} means that recently
accepted or presented papers may not yet be indexed in major databases.
Therefore, the review reflects the state of the literature up to July 2026.

Finally, some studies did not provide complete information about hardware,
software, datasets, or reproducibility, which limited the depth of certain
comparisons.
